%% document class
\documentclass[10pt,oneside,leqno,openright]{memoir}

\input{macros}
\newcommand*{\sto}{\rightsquigarrow}


\begin{document}

\frontmatter
\thispagestyle{empty}
\begin{center}
  \LARGE
  {\bfseries Synthetic Homotopy Theory}\\[1cm]
  \large Exercises and Solutions for the\\[5mm]
  \Large Midlands Graduate School (MGS)\\
  \large 13--17 April 2026, Nottingham, UK\\[1cm]
  by\\
  Axel Ljungström and Aref Mohammadzadeh\\
  School of Computer Science\\
  University of Nottingham\\[1cm]
  Heavily based on notes from MGS 2024 by
  Ulrik Buchholtz and Mark Williams\\
  \Large NB\enspace Subject to change: Refresh if needed
\end{center}

\vfill

%\begin{abstract}
%  These lecture notes aim to introduce the field of synthetic homotopy theory to a (mainly) computer science audience. That is, we shall assume that the reader cares about types,
%\end{abstract}

\clearpage

\tableofcontents*

\mainmatter

%% \setcounter{chapter}{-1}

%% \chapter{Introduction}

\chapter{Truncated and connected types}
%\textcite{Shulman-logic-of-space}


\section{Exercises}
\begin{enumerate}
    \item Show that for $n \geq m$ and any type $A$ we have
    \[ \Trunc{\Trunc{A}_n}_m \simeq \Trunc{A}_m. \]

  \item\label{ex:n-conn-orth}
    Show that a type $X$ is $n$-connected if and only if, for every $n$-type $Y$,
    the diagonal map
    \[ \delta: Y \to (X \to Y), \quad y\mapsto_\delta (x\mapsto y),
    \]
    is an equivalence.

  \item Show that the map $\trunc\blank_n : X \to \Trunc X_n$ is $n$-connected.

  \item Show that a type $X$ is $n+1$-truncated if and only if the diagonal map of $X$,
    \[\delta_X : X \to X\times X,\quad x\mapsto (x,x),\]
    is $n$-truncated.
    
    \item Consider a map $f : X \to Y$.
      We define its \textbf{$n$-image} to be the type $\im_n(f) :=
      (y : Y) \times \Trunc{\fib_f(y)}_n$.
    \begin{enumerate}
    \item Show that we can factorise $f$ through the $n$-image, that is,
      construct a commuting diagram:
        \[\begin{tikzcd}
            X \ar[rr, "f"] \ar[dr, swap, "p"] & & Y \\
            & \im_n(f) \ar[ur, swap, "i"]
        \end{tikzcd}\]
    \item Show that $p$ is $n$-connected and $i$ is $n$-truncated.
    \item Show that if $Z$ is any type, $g : X \to Z$ is $n$-connected, and $h:Z\to Y$ is $n$-truncated with $f = h\circ g$, then there is a unique equivalence $e : \im_n(f) \to Z$ with $e\circ p = g$ and $h\circ e = i$:
      \[
        \begin{tikzcd}
      & Z \ar[dr,"h"] & \\
      X \ar[rr,pos=0.7,"f"]\ar[dr,"p"'] \ar[ur,"g"]& & Y\\
      & \im_n(f)\ar[ur,"i"']\ar[uu,dotted,pos=0.3, "e"] &
    \end{tikzcd}
  \]
  (\emph{Hint})\enspace It's easier if you assume $g$ and $h$ are projections of
  dependent sums.
    \end{enumerate}

    \item (Eckmann--Hilton). All paths will take place in a fixed type $X$.
    \begin{enumerate}
        \item Given paths $p_0, p_1 : x = y$ and $q_0, q_1: y = z$, and higher paths $r : p_0 = p_1$ and $s : q_0 = q_1$, define their \textbf{horizontal multiplication} $s \ast r : q_0 \cdot p_0 = q_1 \cdot p_1$ by path induction on $r$ and $p_0$.
        \[\begin{tikzcd}
    	x & y & z
    	\arrow[""{name=0, anchor=center, inner sep=0}, "{p_1}"', bend right=40, from=1-1, to=1-2]
    	\arrow[""{name=1, anchor=center, inner sep=0}, "{q_1}"', bend right=40, from=1-2, to=1-3]
    	\arrow[""{name=2, anchor=center, inner sep=0}, "p_0", bend left=40, from=1-1, to=1-2]
    	\arrow[""{name=3, anchor=center, inner sep=0}, "q_0", bend left=40, from=1-2, to=1-3]
    	\arrow["s", shorten <=3pt, shorten >=3pt, Rightarrow, from=3, to=1]
    	\arrow["r"', shorten <=3pt, shorten >=3pt, Rightarrow, from=2, to=0]
        \end{tikzcd}\]

        \item Given paths fitting into the following diagram
    	\[\begin{tikzcd}
        x && y && z
        \arrow[""{name=0, anchor=center, inner sep=0}, "{p_1}"', bend right=40, from=1-1, to=1-3]
    	\arrow[""{name=1, anchor=center, inner sep=0}, "{q_1}"', bend right=40, from=1-3, to=1-5]
    	\arrow[""{name=2, anchor=center, inner sep=0}, "{p_0}", bend left=40, from=1-1, to=1-3]
    	\arrow[""{name=3, anchor=center, inner sep=0}, "{q_0}", bend left=40, from=1-3, to=1-5]
    	\arrow[""{name=4, anchor=center, inner sep=0}, from=1-1, to=1-3]
    	\arrow[""{name=5, anchor=center, inner sep=0}, from=1-3, to=1-5]
    	\arrow["{r_0}"', shorten <=2pt, shorten >=2pt, Rightarrow, from=2, to=4]
    	\arrow["{r_1}"', shorten <=2pt, shorten >=2pt, Rightarrow, from=4, to=0]
    	\arrow["{s_0}", shorten <=2pt, shorten >=2pt, Rightarrow, from=3, to=5]
    	\arrow["{s_1}", shorten <=2pt, shorten >=2pt, Rightarrow, from=5, to=1]
        \end{tikzcd}\]
        show we have the \textbf{interchange law}:
        \[ (s_1 \ast r_1) \cdot (s_0 \ast r_0) = (s_1 \cdot s_0) \ast (r_1 \cdot r_0)  \]

        \item Deduce that for $n \geq 2$, the path composition in $\Omega^n(X, x)$ is commutative, and thus that $\pi_n(X, x)$ is abelian.
    \end{enumerate}
    \item (The Fundamental Theorem of Identity Types)
    
    Let $R : X \to X \to \UU$ be a reflexive binary relation. The following are equivalent   \begin{enumerate}
        \item The obvious map\footnote{i.e.\ the map defined by path induction which
          sends $\refl_x$ to the term of type $R(x,x)$ given to us by
          reflexivity of $R$} $x =_A y \to R(x,y)$  is an equivalence.
        \item For any $x :X$, the total space $(y : X) \times (R(x ,y))$ is contractible.
    \end{enumerate}
 \item Recall that $\UU^{\le n} := (X: \UU) \times \isType{n}(X)$ is an $(n+1)$-type. This allows us to define $R : \Trunc{X}_{n+1} \times \Trunc{X}_{n+1} \to \UU^{\le n}$ by the truncation elimination: 
    \[R(\trunc{x}_{n+1} , \trunc{y}_{n+1}) := \Trunc{x = y}_{n}\]
    %

    Apply the fundamental theorem of identity types with the above
    family to deduce that $(\trunc{x}_{n+1} =_{\Trunc{X}_{n+1}} \trunc{y}_{n+1}) \simeq \Trunc{x = y}_n$
    
  \item Show that when $X$ is connected, that $\pi_n(X, x_0) \cong \pi_n(X, x_1)$ are merely isomorphic for any $x_0, x_1 : X$, that is, construct an element of the type
    \[
      (x_0, x_1: X) \to \Trunc{\pi_n(X,x_0) \simeq \pi_n(X,x_1)}.
    \]

    \item We define the observational equality type family of the natural numbers $E : \NN \to \NN \to \UU$  by
        \begin{align*}
            E(0,0) &= 1 \\
            E(\mathrm{succ}(n), 0) &= 0 \\
            E(0, \mathrm{succ}(n)) &= 0 \\
            E(\mathrm{succ}(n), \mathrm{succ}(m)) &= E(n,m)
        \end{align*}
        \begin{enumerate}
            \item Define a natural map $n = m \to E(n,m)$ for all $n, m : \NN$.
            \item Use the fundamental theorem of identity types to show this natural map is an equivalence.
        \end{enumerate}


  \item Show that $\Trunc X_0$ satisfies the universal property of the set
    quotient of $X$ modulo the equivalence relation $\sim$, where $x\sim x'$ if
    and only if $\Trunc{x=_X x'}$.

  %% \item Given two functions $f : X_1 \to Y$ and $g : X_2 \to Y$, we
  %%   can define their pullback $X_1 \times_Y X_2 := (x_1 : X_1) \times
  %%   (x_2 : X_2) \times (f(x_1) = g(x_2))$. Letting $\UU/Y : = (X :
  %%   \UU) \times (X \to Y)$, this can be viewed as a product $\times :
  %%   \UU/Y \times \UU/Y \to \UU/Y$ defined by $(X_1,f) \times (X_2,g)
  %%   := (X_1 \times_Y X_2 , (x_1,x_2,p) \mapsto f(x_1))$
  %%   \begin{enumerate}
  %%   \item Define the functorial action of $\times$. For this question, you should is a triple $(\alpha,\beta,\gamma)$.
  %%   \item An Define the functorial action
  %%   \end{enumerate}


% \begin{enumerate}
%   \item The obvious way of defining the type of arrows over the universe $(x:X)$
% \end{enumerate}
% This means, in particular, that we have an automorpism $(x:X) \times (X \to Y) \simeq (x : X) \times(X \to Y)$ defined by
% \[(X \xrightarrow{f} Y) \mapsto ((y : Y) \times (\fib_f(y)) \xrightarrow{\fst} Y).\]
% As a result, anytime we wish to prove a property ranging over the the slice $\UU/Y := $
\end{enumerate}
\newpage



\raggedright
%\printbibliography

\end{document}

%%% Local Variables:
%%% mode: latex
%%% TeX-master: t
%%% End:
